\documentclass[11pt,a4paper]{article}
\usepackage{../shared/luxcover}
\usepackage{../shared/proposal}
\usepackage[utf8]{inputenc}\usepackage[T1]{fontenc}
\usepackage{amsmath,amssymb}
\usepackage{array}\usepackage{tabularx}
\usepackage{listings}\usepackage{xcolor}
\input{../shared/lstlang}

\newcommand{\code}[1]{\texttt{\small #1}}
\newcommand{\pv}[1]{{\color{muted}\sffamily\scriptsize[#1]}}
\newcolumntype{Y}{>{\raggedright\arraybackslash}X}
\newcommand{\cond}{$^{\dagger}$}

\title{Scoring the Migration}
\author{Zach Kelling\\\textit{Lux Industries, Inc.}\\\texttt{research@lux.network}}
\date{2026}

\begin{document}
\luxcoverpage
\sloppy

\noindent{\color{muted}\footnotesize\sffamily RUBRIC APPLIED \quad\textbullet\quad Morpheus after migration \quad\textbullet\quad Hanzo AI \& Lux Industries \quad\textbullet\quad 2026}\par
\vspace{0.8ex}
\noindent{\LARGE\bfseries Scoring the Migration}\par
\vspace{0.4ex}
\noindent{\itshape Does the proposed Lux end-state distribute the powers, or rename the multisig. The rubric is applied to our own proposal, and to our own stack, on the same terms as Morpheus.}\par
\vspace{0.7ex}
\noindent{\color{rule}\rule{\linewidth}{1.1pt}}\par
\vspace{1.0ex}

\section{What is being scored}

The target is Morpheus as the adoption series proposes to leave it: upgrade
authority behind \code{lux/dao} (a Safe with a Governor module and a timelock
delay), a vote-escrow token \code{vMOR}, settlement moved to \code{mord} on a Lux
L2, with the L2 anchoring to Lux primary-network finality. The question is not
whether that is better than the anonymous 5-of-9 Morpheus runs today. It is
whether the end-state actually distributes the six powers, or whether it moves
them from a multisig to a stack that one party still controls under better names.
This is graded from the code the proposal names, not from the proposal's prose,
and the same standard is turned on the Lux contracts and nodes that would host it.

Scores are read best-faith: where a power depends on a migration step (transfer
\code{owner()} to the timelock, remove the Safe's direct signers, open the
validator set), the score assumes the step is completed and is marked
conditional\cond. The heal-thyself section then reports which of those steps our
own live deployments have not in fact taken, because a conditional score the
sponsor never satisfies is theater by the rubric's own definition.

\section{The vector}

\begin{table}[h]\footnotesize\setlength{\tabcolsep}{5pt}\renewcommand{\arraystretch}{1.2}
\begin{tabularx}{\linewidth}{@{}l c Y Y@{}}
\toprule
\textbf{Power} & \textbf{Score} & \textbf{Measured state (file:line)} & \textbf{What the score depends on}\\
\midrule
Upgrade & 1\cond & \code{owner()} moves to a delay; but the ``timelock'' is a module parameter the owner tunes (\code{ModuleGovernorV1.sol:404} \code{updateTimelockPeriod onlyOwner}) and the module is UUPS-upgradeable by its owner (\code{:200}) & Safe signers removed and module owned by the DAO, not by a multisig that can still sign directly\\
Custody & 1\cond & \code{MOR} mint stays uncapped, now behind the delay (\code{MOROFT.sol:50,61}, unchanged by the migration); session escrow returns on close; slashing is app-level in \code{mord}, not primary-level & A supply cap, and the delay actually gating the owner-minter path\\
Halt / censor & \textbf{0} & An L2 from \code{l2-template} produces blocks from one operator's five pods (\code{init.sh} \code{node.yaml replicas:5}, one namespace; MPC and bridge single-replica); ADOPT-01 states ``Validators are Lux-managed'' & Opening the L2 producer set to independent operators; a forced-inclusion / escape path\\
Governance & 1\cond & Governor to timelock to execution exists, but \code{vMOR} weight is flat 1:1 (\code{VotesERC20StakedV1.sol} \code{stake} to \code{\_mint}), quorum is 4\% \pv{zoo.json}, and 76\% of emission is owner-directed (\code{DistributionV6.sol:175} \code{onlyOwner}) & A vote the emission-holder cannot alone carry; Safe removed so execution is vote-gated\\
Exit & 1 & Code is open and forkable; \code{MOR} is self-custodied on Ethereum/Arbitrum/Base; funds bridged onto the single-operator L2 have no withdrawal independent of that operator & An escape hatch from the L2 that does not require the operator to be live\\
Liveness & \textbf{0} & The compute L2 halts if the one operator's cluster, MPC signer, or bridge stops; the base-chain \code{MOR} contracts are unaffected & Many independent validators/RPCs for the chain that carries settlement\\
\bottomrule
\end{tabularx}
\end{table}

\noindent\pv{\cond conditional: assumes a migration step our live reference
deployments have not performed; see heal-thyself.}

\vspace{0.4ex}
\noindent Vector \([\,1^\dagger,\ 1^\dagger,\ 0,\ 1^\dagger,\ 1,\ 0\,]\). The
real-decentralization ceiling is the lowest score among load-bearing powers. Halt
and Liveness are 0, so the ceiling of the migrated compute chain is \textbf{0},
the same ceiling as Morpheus today. The migration's genuine gains sit on the
authority axes (Upgrade, Custody, Governance move from 0 toward 1), and every one
of those gains is conditional on a step the sponsor has not yet taken in its own
production stack.

\section{The two zeros set the ceiling}

The proposal's strongest security claim is inherited post-quantum finality from
Quasar. That claim is about finality, which is the order and irreversibility of
blocks once produced. It is not about who produces them. An L2 stood up by
\code{l2-template} is a Kubernetes \code{StatefulSet} of five node replicas in one
operator's namespace, beside a single MPC signer and a single bridge
\pv{verified: \code{l2-template/init.sh}}. ADOPT-01 is explicit that ``Validators
are Lux-managed'' \pv{verified: adopt-01}. Both readings fail the client's own
test: if Lux runs the validators it is one org, and if Morpheus runs the five pods
it is one org. The permissionless staking of the Lux \emph{primary} network
(anyone meeting \code{MinValidatorStake} may validate via
\code{AddPermissionlessValidatorTx}, \pv{verified: \code{node/vms/platformvm}})
does not transfer to the L2's producer set; the L2 recognizes primary-network
finality but sequences its own blocks. So Halt/censor is 0: one operator produces,
a user can be excluded, and there is no forced inclusion. Liveness is 0 for the
same reason: the chain that carries settlement stops when that one operator stops.

This is not worse than Morpheus today only because Morpheus today runs on Base,
which is itself a single-sequencer rollup. The migration does not fix the axis; it
moves the single sequencer from Coinbase's infrastructure to one org's five pods.
Two independent facts compound it: Lux primary staking has no slashing of
principal (\code{grep -ri slash node/vms/platformvm node/snow} is empty
\pv{verified}), so even the finality layer's validators are bonded but not
slashable; and the realized stake distribution, the number that would set the
primary network's Nakamoto coefficient in practice, is not something the
repository establishes \pv{unmeasured: no on-chain census in-tree}. We do not
assert a coefficient we have not counted.

\section{Governance is binding but capturable}

Take the migration as completed: \code{owner()} sits behind the timelock, the Safe
no longer signs directly, and every mint, upgrade, and pool edit is propose, vote,
delay, execute. That is a real improvement over an off-chain forum, and it earns
the 1. It does not earn the 2, and the reason is the specific mechanism, not a
distribution that ``might'' thin out.

\code{vMOR} is named for vote-escrow and the token-design note describes weight
``scaled by lock length,'' so that ``a long holder and a mercenary one weigh
differently'' \pv{proposal: mor-token-design.md}. The contract the proposal names,
\code{VotesERC20StakedV1}, does not do this. \code{stake(amount)} mints voting
power one-to-one with the amount staked (\code{\_mint(msg.sender, amount\_)}), and
the only time constraint is a single global \code{minimumStakingPeriod} checked at
\code{unstake} \pv{verified: repo}. There is no per-lock duration and no
duration-weighting anywhere in the contract set \pv{verified: no lock-length term
in \code{dao/contracts}}. Vote weight is therefore proportional to \code{MOR}
staked, full stop.

That matters because of where \code{MOR} comes from. Seventy-six percent of
emission is assigned by the owner through \code{manageUsersInPrivatePool}, an
\code{onlyOwner} call (\code{DistributionV6.sol:175} \pv{verified}). Quorum in the
live config is 4\% \pv{verified: \code{zoo.json}}. A vote whose weight is
proportional to a balance the emission controller allocates, at a 4\% quorum, is a
vote the emission controller can meet alone. The capture vector the design says it
closes is left open by the contract the design cites. Under the rubric's own
Governance test, ``insiders alone meet quorum'' is the theater case; the timelock
and the on-chain execution are what keep this at 1 rather than 0, but the 1-to-2
step, a vote no insider can carry, is not delivered and is undermined by the flat
weighting.

\section{Upgrade and custody: real progress, still conditional}

The authority axes are where the migration does honest work, and it is worth
stating plainly. Today Morpheus has no execution delay anywhere in either contract
tree and no vote surface at all; the compute diamond is an \code{OwnableDiamond}
(\code{LumerinDiamond.sol:9}) whose owner can recut any facet in one transaction,
and \code{MOROFT} mint is uncapped behind an owner-set minter (\code{:50,61})
\pv{verified}. Putting \code{owner()} behind a timelocked Governor converts each of
those one-key actions into a delayed, visible, on-chain process, which is the
0-to-1 move and a genuine one.

The conditions are the whole story. Upgrade reaches 1 only if the Safe's direct
signers are removed, because the \code{lux/dao} timelock is not a standalone
\code{TimelockController}: it is the Azorius module's internal \code{timelockPeriod},
and both that period (\code{updateTimelockPeriod}, \code{onlyOwner}) and the module
itself (UUPS \code{\_authorizeUpgrade}, \code{onlyOwner}) answer to the module
owner (\code{ModuleGovernorV1.sol:404,200} \pv{verified}). If the owner is a Safe
whose signers can still sign, the delay is a setting that same group can zero.
Custody reaches 1, not 2, because the migration explicitly does not change
\code{MOROFT}; the uncapped owner-minter path survives, now behind the delay
rather than removed or capped. \code{mord} adds slashing, but at the application
layer for providers, not to the validators that secure the chain.

\section{Heal-thyself: the conditions our own stack has not met}

The conditional scores above assume steps our production deployments have not
taken. Disclosing this is the point: selling ``real decentralization'' while our
own chains retain admin authority would be the exact hypocrisy the client is
trying to avoid.

\begin{description}
\item[Our production timelock keeps its deployer as admin.] The Lux governance
deploy constructs \code{Timelock(1 days, [deployer], [address(0)], deployer)}
(\code{standard/contracts/script/DeployMultiNetwork.s.sol:254} \pv{verified}). The
deployer is the sole proposer and the retained admin role, and the script never
renounces it (\code{grep} for \code{renounce}/\code{revokeRole} is empty).
\code{vLUX} and the \code{GaugeController} are deployed but not wired as proposers.
As scripted, that is a one-day delay a single EOA administers and alone can
propose into: the governance contracts exist, the vote does not gate execution.
This is the same finding we level at Morpheus, in our own tree. \pv{verified:
deploy script; on-chain role state not independently re-confirmed here}
\item[Our launch multisig is one seed.] \code{governance.md} specifies all five
launch signers derived from a single \code{LUX\_MNEMONIC}, with rotation to
independent hardware keys planned before a deadline (2026-08-08) that has now
passed, status ``TBD'' \pv{verified}. Until rotation, the effective control set of
the launch Safe is one, which is the Ronin shape the rubric names. The doc is
candid that this is time-bounded bringup, which is honesty, not theater; but the
rotation is a claim about the future, not a measured state.
\item[Our verified settlement is not yet verified.] We would offer \code{mord} and
\code{lux/ai} as the un-fakeable base. The current \code{lux/ai} settlement path,
\code{aivm.go} \code{SubmitResult}, builds a reward receipt from a self-reported
result (provider id, compute time, output) and calls \code{SubmitReceipt}
(\code{:311}) with no signature, countersignature, attestation, or quorum check.
The attestation \code{verifier} runs only at provider registration (\code{:216,223}),
and \code{SettleQuorum} (\code{rewards.go:407}) is off this path; the miner polls
the public \code{/v1/bc/A/ai/pendingTasks} and is paid the self-reported reward
(\code{miner.go}) \pv{verified}. The L2 strategy note already flags this and orders
it first \pv{disclosed: morpheus-as-lux-l2.md}. It is disclosed, and it is not yet
fixed.
\item[Our second governance tenant is a front end.] ADOPT-03 cites \code{zoo/vote}
as the stack ``white-labeled through one config,'' implying a live instance. Its
\code{config/zoo.json} has every contract address null: \code{safeFactory},
\code{moduleGovernor}, \code{governor} \pv{verified}. It is a branded UI pointing at
no deployed contracts, not a running governance deployment.
\end{description}

\section{Honesty flag and verdict}

\noindent\textbf{Honesty flag: mixed.} The architecture-level trade-offs are
disclosed. The notes state that the L2 surrenders producer-set control to a
Lux-managed set, that the settlement bypass must be fixed first, and that a raise
clears only at real demand \pv{disclosed}. Those are immaturity, not theater.
Three claims are not disclosed and are contradicted by the code: \code{vMOR}
presented as duration-weighted when the named contract is flat 1:1; ``MOR becomes
credibly decentralized'' and ``real governance'' presented as delivered when our
own reference stacks retain a deployer-admin timelock and an operational Safe; and
\code{zoo/vote} presented as a live tenant when its addresses are null. A property
claimed but not present in the code is the gap the rubric exists to measure.

\noindent\textbf{Verdict: theater on the load-bearing powers, genuine but
conditional progress on the rest.} The mechanism exists. It does not, as
specified, distribute the powers that carry funds and liveness. Halt and Liveness
stay at 0 because settlement moves onto a single-operator L2. Governance is
binding but capturable by the party that already directs 76\% of emission, because
the vote is weighted by a balance that party allocates. Upgrade and Custody
improve from 0 toward 1, conditional on removing an admin our live deployments
still hold. The net effect of the migration as drawn is to replace an anonymous
5-of-9 with a named stack that a single operator and a single emission-holder can
still control: centralized under a nicer name, which is precisely the outcome the
client asked us to avoid.

Three changes make it real, and none is delivered by the default path.
\emph{Open the L2 producer set} to independently operated validators, or stop
describing finality-anchoring as validator decentralization. \emph{Make the vote
capture-resistant}: weight it by committed lock duration as the design claims, and
set quorum above the owner-directed emission share, or the emission controller is
the government. \emph{Actually renounce the admin}: remove the Safe's direct
signers and renounce the timelock's deployer-admin role, on Lux first, so the
delay we sell is one we ourselves cannot bypass. Until those land, the migration is
a better starting point than Morpheus has today, honestly graded as a 1-ceiling
proposal with two load-bearing zeros, not the finished decentralization it reads as.

\vspace{0.8ex}
\noindent{\color{muted}\sffamily\footnotesize RUBRIC APPLIED. Method in ``Real
Decentralization''; governance path in ADOPT-03; settlement in ADOPT-02; the L2
trade in ADOPT-01. \; \textbullet\; \texttt{research@lux.network}}

\end{document}
